\documentclass[12pt,a4paper]{report}

\usepackage[T1]{fontenc}
\usepackage[utf8]{inputenc}
\usepackage[hungarian]{babel}
\usepackage{lmodern}
\usepackage{tikz}
\usetikzlibrary{shapes,arrows,positioning}
\usepackage{graphicx}
\usepackage{tabularx}
\usepackage{booktabs}
\usepackage{listings}
\usepackage{xcolor}
\usepackage{pdfpages}
\usepackage{hyperref}


\lstset{
    basicstyle=\ttfamily\small,
    columns=fullflexible,
    breaklines=true,
    frame=single,
    showstringspaces=false,
    literate={á}{{\'a}}1 {é}{{\'e}}1 {í}{{\'i}}1 {ó}{{\'o}}1 {ö}{{\"o}}1
             {ő}{{\H{o}}}1 {ú}{{\'u}}1 {ü}{{\"u}}1 {ű}{{\H{u}}}1
             {Á}{{\'A}}1 {É}{{\'E}}1 {Ö}{{\"O}}1 {Ü}{{\"U}}1
}

\title{Interaktív webalkalmazás a Számításelmélet alapjai I–II tárgyhoz}
\author{Szabó Anna}
\date{2026}

\begin{document}
\maketitle
\includepdf{temabejelento.pdf}
\tableofcontents

\part{Bevezetés}

\chapter{Motiváció és problémafelvetés}

Az egyetemi évek alatt számos olyan tárgyat tanulunk, amelyek megértése
jelentős absztrakciós képességet igényel, és sokszor nagyon elvontnak,
nehezen megfoghatónak tűnnek. Hallgatóként én is többször találkoztam olyan
tananyaggal, ahol a fogalmak és folyamatok elsőre nehezen voltak értelmezhetők.
Különösen igaz ez azokra a témákra, amelyekben egy rendszer működésének megértéséhez
több egymásra épülő állapotot vagy lépést kell végigkövetni.
Az ilyen absztrakt gondolatmenetek papíron gyakran nehezen követhetők,
hiszen egy folyamat vagy modell működését látni egészen más élmény,
mint pusztán olvasni róla.

Erre jó példa a Számításelmélet alapjai I--II. tárgy,
amely olyan témaköröket dolgoz fel, mint a formális nyelvek,
a grammatikák, a véges automaták, a reguláris kifejezések,
a veremautomaták vagy a Turing-gépek.
Ezek a fogalmak a hallgatók számára gyakran nehezen értelmezhetők, mivel
működésük és viselkedésük hagyományos tananyagok segítségével nem mindig
szemléltethető egyszerűen.
A tananyag elsajátításához általában előadásanyagok, prezentációk és
órai jegyzetek állnak rendelkezésünkre. Bár ezek fontos szerepet
töltenek be az oktatásban, és ezek adják a tudásunk alapját, a
mélyebb megértés sokszor a gyakorlatokra, illetve az önálló otthoni
tanulásra hárul.

Saját tapasztalatom szerint ezeknél a témáknál a megértés kulcsa
éppen az, hogy az ember minden állapotot és lépést végigkövet. Papíron
ez azonban lassú: gyakran minden állapotot külön le kell rajzolni, egy
elrontott lépés után pedig elölről kell kezdeni. Ezért az ilyen
témáknál mi, hallgatók gyakran nagyon sok időt töltünk az anyag önálló
feldolgozásával.

A szakdolgozat erre a problémára kíván megoldást nyújtani egy olyan
webalapú oktatási rendszer fejlesztésével, amely a tantárgy
legfontosabb témaköreit interaktív, vizualizációra épülő formában
dolgozza fel. A rendszer a témakörökhöz elméleti összefoglalót,
lépésenként követhető szimulációkat (például determinisztikus véges
automata, grammatikai levezetés) és önellenőrző feladatokat kínál,
így a lépésenkénti követés gyorssá és hibatűrővé válik.
% TODO: a felsorolt példákat a végén egyeztetni az elkészült modulokkal

\chapter{Célkitűzések}
A szakdolgozat célja egy interaktív és vizualizációs oktatási webalkalmazás
tervezése és fejlesztése, amely támogatja a hallgatókat a Számításelmélet
alapjai I--II. tantárgy témaköreinek elsajátításában.
A fejlesztés során olyan rendszer létrehozása a cél,
amely egyetlen felületen egyesíti az elméleti háttéranyagot, a vizuális
szemléltetést és az interaktív gyakorlási lehetőségeket. A hallgatók így
nem csupán olvashatnak az egyes modellekről és algoritmusokról, hanem
működésüket lépésről lépésre követhetik, és saját bemenetek segítségével
ki is próbálhatják.
A rendszer fontos célja továbbá, hogy támogassa az önálló tanulást,
valamint segítséget nyújtson a gyakorlásban és a vizsgákra való
felkészülés során.
A kitűzött célok az alábbiakban foglalhatók össze:

\begin{itemize}
    \item a Számításelmélet alapjai I--II. főbb témaköreinek feldolgozása;
    \item elméleti összefoglaló a tárgy jelöléseivel, matematikai képletekkel;
    \item saját felhasználói bemenetre működő, lépésenként követhető szimulációk;
    \item azonnali visszajelzés és önellenőrzés: kvízek, feladatok;
    \item böngészőből, telepítés nélkül elérhető felület;
    \item bővíthető, moduláris alkalmazásarchitektúra kialakítása.
\end{itemize}

\chapter{Hatókör és korlátok}

\section{Hatókör}

% TODO: a számokat és a modullistát a végén egyeztetni a kóddal
% (packages/content/src/topics, apps/web/src/modules/index.tsx)
A szakdolgozat egy webalapú, egyfelhasználós oktatási alkalmazás
tervezését és megvalósítását foglalja magában, amely a Számításelmélet
alapjai I--II. tárgy témakörét dolgozza fel Minden témakörhöz elméleti összefoglaló
tartozik, amelyet több helyen interaktív modul egészít ki.

rendszer összesen x témakört dolgoz fel, amelyek xx nagyobb csoportba sorolhatók: 
alapfogalmak, reguláris nyelvek, környezetfüggetlen nyelvek és környezetfüggő nyelvek...
 Az egyes témakörökhöz elméleti összefoglalók, valamint interaktív modulok tartoznak, 
 amelyek lehetővé teszik a modellek és algoritmusok működésének lépésről lépésre történő megfigyelését.
 
A megvalósított modulok közé tartozik többek között a DVA-szimulátor, a grammatikalevezető, a Chomsky-osztályozó, a 3-as normálforma vizualizáló, a reguláris kifejezés ellenőrző, a nyelvfelfedező, a szójáték és a kvízrendszer.

% TODO: NDVA-szimulátor elkészült? Ha igen, kerüljön a fenti listába
% TODO: ha a témakör-táblázat bekerül a fejlesztői dokumentációba, ide hivatkozás:
% A témakörök és modulok részletes áttekintését a \ref{tab:temakorok}.~táblázat tartalmazza.

A dolgozat része a fejlesztői dokumentáció, vagyis
a rendszer architektúrájának bemutatása, a tervezés
dokumentálása, az implementáció ismertetése, valamint a megvalósított
funkciók tesztelése és értékelése. Végül egy felhasználói dokumentáció
is készül, amely a rendszer használatát segíti.

\section{Korlátok}

A dolgozat keretében megvalósított rendszer oktatási célú
prototípusnak tekinthető. Bár alkalmas a tanulás támogatására, nem
célja egy teljes körű e-learning vagy tanulmányi rendszer kiváltása.
A rendelkezésre álló idő és fejlesztési kapacitás miatt a tantárgy teljes
 tananyagának részletes feldolgozása nem valósul meg. 
 A rendszer elsősorban azokra a témakörökre összpontosít, amelyek különösen jól támogathatók vizualizációval
  és interaktív szimulációkkal.
A Számításelmélet alapjai II. tárgy csak részben kerül feldolgozásra.
A szimulációk kisméretű, véges bemenetekre lettek tervezve.
Nincs felhasználókezelés, a haladás nem tárolódik központilag.

\section{Feltételezések}

A rendszer használata modern webböngésző meglétét feltételezi. A
felhasználók rendelkeznek alapvető informatikai ismeretekkel, valamint
a Számításelmélet alapjai I--II. tárgy követéséhez szükséges
előismeretekkel. Mivel a tananyag az előadásokra épül, azokat nem
helyettesíti, hanem kiegészíti. A tervezés során feltételezem, hogy a webes
technológiák és a kliensoldali vizualizációs megoldások elegendő
teljesítményt biztosítanak az interaktív modulok futtatásához.

\chapter{A dolgozat fő eredményei}

\begin{itemize}
    \item egy interaktív, webalapú oktatási rendszer tervezése és
          megvalósítása a Számításelmélet alapjai I--II. tárgy támogatására;
    \item olyan interaktív vizualizációs és szimulációs eszközök fejlesztése, amelyek
     lehetővé teszik a hallgatók számára a számításelméleti modellek működésének lépésről 
     lépésre történő követését;
    \item a tananyag egységes, típusos, blokkalapú tartalommodellbe
          szervezése;
    \item moduláris és bővíthető szoftverarchitektúra kialakítása;
    \item interaktív szimulációs és kiértékelő mechanizmusok fejlesztése;
    \item a rendszer tervezésének, implementációjának és tesztelésének
          dokumentálása.
\end{itemize}

A dolgozat eredménye egy olyan oktatási segédeszköz, amely a
hagyományos tananyagokat kiegészítve támogatja az önálló tanulást, a
szemléletesebb megértést és a gyakorlati tapasztalatszerzést.

\chapter{A dolgozat felépítése}
A dolgozat két fő részre tagolódik.
A felhasználói dokumentáció a rendszer felépítését és használatát
ismerteti a végfelhasználó szemszögéből. Ez kiterjed az egyes modulok
bemutatására és az interaktív, szimulációs eszközök használatára
példákon keresztül.

A fejlesztői dokumentáció a műszaki hátteret mutatja be. Itt kap helyet
a követelményelemzés, a tervezési döntések indoklása, az architektúra
ismertetése és az implementáció bemutatása. Ezenkívül ez a rész tér ki
a tesztelési folyamatokra, valamint a jövőbeni fejlesztési irányokra
és bővítési lehetőségekre is.
A dolgozatot egy rövid értékelés zárja.

%=========================================
\part{Felhasználói dokumentáció}
%======================================================================

% ha már nagyjából kész!
TODO

%======================================================================
\part{Fejlesztői dokumentáció}
%======================================================================

\chapter{Követelményelemzés és specifikáció}

A fejlesztés és kód írás megkezdése előtt fontos egy rövid felmérést, kutatást
végeznünk és bizonyos szempontok alapján meg kell vizsgálnunk a feladatot.
A rendszer fejlesztésének első lépése a követelmények meghatározása és dokumentálása.
A specifikáció elkészítése.
 A követelményelemzés során tisztázzuk felhasználói igényeket, a rendszerrel szemben támasztott elvárásainkat,
valamint azokat a korlátokat és feltételeket, amelyek a fejlesztést befolyásolják.
 
   
A fejezet a specifikáció és követelménymenedzsment során alkalmazott
lépéseket követi. Ennek megfelelően először a rendszer megvalósíthatóságát
  vizsgálom, majd a követelmények feltárását és elemzését végzem el. 
Ezután bemutatom a funkcionális és nem funkcionális követelményeket, 
ismertetem a rendszer alap modelljeit, végül pedig felvázolom azokat a feltételeket, amelyeket,
ha a rendszer teljesít, akkor elfogadható.

%----------------------------------------------------------------------
\section{Megvalósíthatósági elemzés}
%----------------------------------------------------------------------
% ea02, 4. dia: erőforrások, költség/időtartam, üzemeltetés/karbantartás
% - humán: egy fejlesztő (én), a tárgyat elvégzett hallgatóként ismerem az anyagot
% - hardver: nincs szükség szerverre, minden a böngészőben fut
% - szoftver: csak ingyenes, nyílt forrású eszközök (TypeScript, React, Vite, pnpm)
% - technikai: a számítások kis, korlátos bemenetekre készülnek
%   (shortlex felsorolás, ellenpélda-keresés hosszkorlátos) -> böngészőben is azonnaliak
% - időtartam: egy félév; ami nem fér bele -> Prioritizálás szakasz
% - üzemeltetés: statikus GitHub Pages, ingyenes; karbantartás = a tananyag TS-fájljainak frissítése
Itt vizsgálom a rednszer, a szoftver és a fejlesztés megvalósíthatóságát,
 figyelembe véve a rendelkezésre álló erőforrásokat, a költségeket és a rendelkezésre álló időtartamot,
  valamint az üzemeltetés és karbantartás kéréseit.

A rendszer fejlesztését egyetlen fejlesztő végzi, nincs szükség külső erőforrás bevonására.
 A fejlesztés közvetlen költsége gyakorlatilag nulla, mivel a projekt kizárólag ingyenesen elérhető,
nyílt forráskódú technológiákra épül.
Hardverigény szempontjából a rendszer kedvező helyzetben van, mivel minden számítás a felhasználó
böngészőjében fut le, így nincs szükség külön szerver infrastruktúrára vagy felhőszolgáltatásra.
Az alkalmazás statikus weboldalként publikálható, például GitHub Pages szolgáltatáson keresztül,
amely ingyenes üzemeltetést tesz lehetővé.

A rendszer által végzett algoritmusok viszonylag kis méretű véges bemeneteken dolgoznak, ez által
 azok átlagos felhasználói eszközökön is gyorsan végrehajthatók.

A fejlesztés időkeretét egy egyetemi félév adja. Ennek megfelelően a funkciók prioritások szerint 
kerülnek megvalósításra, és az alacsonyabb prioritású elemek szükség esetén későbbi bővítésként kezelhetők.

A karbantartása egyszerű, ezt támogatja a, moduláris felépítés. 
Az üzleti logika  UI elkülönítése, React által biztosított komponensalapú architektúra. 
Így új funkciók hozzáadása vagy meglévő funkciók módosítása nem igényli a teljes rendszer
átírását. A karbantartást tovább segíti a Git alapú verziókezelés.

%----------------------------------------------------------------------
\section{Követelmények feltárása}
%----------------------------------------------------------------------

\subsection{Források és célcsoport}
A oldal tartalma, az előadások anyagára, órai jegyzetekre, valamit saját hallgatói tapasztalatra épül elsősorban.

 A célcsoport a Számításelmélet alapjai I--II. tárgy hallgatói, vizsgára készülők, 
 valamit az anyag mélyebb megértésére törekvő hallgatók.
% előadásdiák (a tartalommodell ezek elemzéséből készült),
% saját hallgatói tapasztalat; TODO: kérdeztél évfolyamtársakat? ha igen, ide

\subsection{Hatókör}
% thesis.tex: "A rendszer célja és hatóköre" szöveg + tab:in_scope
% JAVÍTANDÓ: nincs grafikus automataszerkesztő; a táblázat oszlopdefiníciója {l X} legyen
% + ide jön a témakör–modul táblázat (tab:temakorok)
A szakdolgozat hatókörébe egy webalapú, oktatást támogató alkalmazás
tervezése és megvalósítása tartozik, amely a Számításelmélet
alapjai I--II. tárgy egyes témaköreinek megértését segíti interaktív
és vizuális eszközökkel. A~\ref{tab:in_scope}.~táblázat a hatókörbe
tartozó funkciókat foglalja össze.

\begin{table}[htbp] 
      \centering 
      \begin{tabularx}{\textwidth}{l X} 
            \toprule 
            \textbf{Funkció} & \textbf{Leírás} \\ 
            \midrule 
            Elméleti tartalomkezelés & 
                  Témakörönként strukturált, blokkalapú tananyag a Számításelmélet 
                  alapjai I--II. tananyagából: definíciók, tételek, lemmák, bizonyítások,
                  példák, képletek, táblázatok és figyelmeztetések stb. \\ 
                  \addlinespace
            Interaktív algoritmusok &
                  A hallgató különböző bemenetek megadásával futtathatja a tananyagban szereplő algoritmusokat. 
                  A rendszer a végrehajtás eredményét és köztes állapotait is megjeleníti. \\
                  \addlinespace
            Folyamatok vizualizációja &
                  Az alkalmazás grafikus elemek segítségével szemlélteti a számításelméleti folyamatokat, 
                  ezáltal támogatva az absztrakt fogalmak megértését. \\
                  \addlinespace
            Önellenőrzés &
                  A rendszer gyakorló feladatokat és kvízeket biztosít az elsajátított ismeretek ellenőrzésére. \\
            Egységes tanulási környezet &
                  A rendszer egy helyen fogja össze a tananyagot, a vizualizációkat és az interaktív eszközöket, 
                  támogatva az önálló felkészülést. \\
                  \addlinespace
            \bottomrule
      \end{tabularx}
      \caption{A rendszer hatókörébe tartozó fő funkciók}
      \label{tab:in_scope}
\end{table}


\subsection{Hatókörön kívül eső elemek}
% thesis.tex: tab:out_of_scope (változatlanul átvehető)

\subsection{Prioritizálás}
% ea02, 5. dia: fontossági sorrend, megvalósíthatatlan követelmények elhagyása
% - kötelező: témakörök + elméleti anyag, DVA-szimulátor
% - fontos: grammatika-levezetés, Chomsky-osztályozó, 3NF, regex-ellenőrző, kvíz
% - ha belefér: NDVA-szimulátor, veremautomata
% - nem készül el (indoklással): Turing-gép, logikai modul -> eltérés a témabejelentőtől
% TODO: a témavezetővel egyeztetni

%----------------------------------------------------------------------
\section{Funkcionális követelmények}
%----------------------------------------------------------------------
% thesis.tex: FK-1 ... FK-10, csak ami teljesül
% + használati eset diagram (thesis.tex: "Használati esetek" szakasz, fig:usecase)

\subsection{Felhasználói történetek}
% ea02, 36-39. dia: given-when-then; a többi modulhoz is hasonlót


%----------------------------------------------------------------------
\section{Nem funkcionális követelmények}
%----------------------------------------------------------------------
% ea02, 7-11. dia kategóriái szerint; a thesis.tex NFK-1 ... NFK-8 ide sorolandó

\subsection{Termékkövetelmények}
% - hatékonyság: azonnali válasz; hosszkorlátos felsorolás/ellenpélda-keresés
% - megbízhatóság: hibás bemenetre hibaüzenet (pl. nem ábécébeli betű, hibás regex)
% - biztonság: nem tárol felhasználói adatot, ezért nem szempont
% - hordozhatóság: modern böngésző, telepítés nélkül
% - felhasználhatóság: kézikönyv nélkül, magyarul, a tárgy jelöléseivel
% TODO: reszponzivitás (NFK-3) ellenőrizni

\subsection{Menedzselési követelmények}
% - környezeti: böngésző, szerver nélkül; KaTeX a képletekhez
%   (thesis.tex "Külső interfészek" táblázata ide olvad be)
% - működési: egy felhasználó; félév közben és vizsgaidőszakban
% - fejlesztési: TypeScript, React, pnpm monorepo, Git

\subsection{Külső követelmények}
% - nincs adatgyűjtés -> adatvédelmi kérdés nem merül fel
% - a tananyag az előadásokra épül -> a forrás megjelölése

%----------------------------------------------------------------------
\section{Rendszermodellek}
%----------------------------------------------------------------------
% ea02, 6. és 42. dia: felület terv, adatformátum leírás

\subsection{A felhasználói felület modellje}
% thesis.tex: "Felhasználói felület modell" (fig:homepage, fig:module)
% TODO: hiányzó képek: img.png, img_1.png

\subsection{Adatformátum: a tartalommodell}
\label{sec:data}
% thesis.tex: "Adatkövetelmények: a tartalommodell" teljes szakasza
% JAVÍTANDÓ: blokktípus-táblázatba warning, funfact, proof; "14 témakör váza" -> 16, kitöltve

%----------------------------------------------------------------------
\section{A követelmények validációja}
%----------------------------------------------------------------------
% ea02, 40. dia: prototípus + tesztesetek
% - a fejlesztés iteratív prototípusokban haladt
% - a felhasználói történetekből lesznek a funkcionális tesztesetek -> Tesztelés fejezet

%======================================================================
% Következő fejezet: Tervezés
% ide megy át a thesis.tex "Áttekintés" (3 csomag + modulregiszter)
% és a "Rendszerarchitektúra" fejezet
%======================================================================

\end{document}
